\documentclass[11pt,a4paper]{article}
\usepackage{amsmath,amssymb}
\usepackage{listings}
\usepackage{hyperref}
\usepackage[margin=1in]{geometry}
\usepackage{fancyhdr}

\pagestyle{fancy}
\fancyhf{}
\fancyhead[L]{Module 06: Pointers}
\fancyhead[R]{C Programming Course}
\fancyfoot[C]{\thepage}

\lstset{
  language=C,
  basicstyle=\ttfamily\small,
  keywordstyle=\bfseries,
  commentstyle=\itshape,
  numbers=left,
  numberstyle=\tiny,
  frame=single,
  breaklines=true,
  tabsize=4
}

\title{Module 06: Pointers}
\author{C Programming Course}
\date{}

\begin{document}
\maketitle
\thispagestyle{fancy}

\section{Introduction to Pointers}
A pointer stores the memory address of another variable.

\begin{lstlisting}
#include <stdio.h>

int main(void) {
    int x = 42;
    int *p = &x;
    printf("Value: %d, Address: %p\n", *p, (void *)p);
    return 0;
}
\end{lstlisting}

\begin{itemize}
  \item \texttt{\&x} -- address-of operator, yields the address of \texttt{x}.
  \item \texttt{*p} -- dereference operator, accesses the value at the address.
\end{itemize}

\section{Pointer Arithmetic}
Pointer arithmetic advances by the size of the pointed-to type.

\begin{lstlisting}
int arr[4] = {10, 20, 30, 40};
int *p = arr;
printf("%d\n", *(p + 2)); /* prints 30 */
\end{lstlisting}

\section{Pointers and Arrays}
An array name decays to a pointer to its first element when passed to a function.

\begin{lstlisting}
#include <stdio.h>

void print_array(const int *arr, int n) {
    for (int i = 0; i < n; i++) {
        printf("%d ", arr[i]);
    }
    printf("\n");
}

int main(void) {
    int data[] = {1, 2, 3, 4, 5};
    print_array(data, 5);
    return 0;
}
\end{lstlisting}

\section{Pointer to Pointer}
\begin{lstlisting}
#include <stdio.h>

int main(void) {
    int val = 100;
    int *p = &val;
    int **pp = &p;
    printf("val = %d\n", **pp);
    return 0;
}
\end{lstlisting}

\section{Pointers and Functions}
Pass pointers to allow a function to modify the caller's data.

\begin{lstlisting}
#include <stdio.h>

void increment(int *n) {
    (*n)++;
}

int main(void) {
    int x = 5;
    increment(&x);
    printf("x = %d\n", x); /* prints 6 */
    return 0;
}
\end{lstlisting}

\end{document}
